% Back matter -- translated by Claude (2026), continuing the voice and
% register of the Burkett-Hutcheson chapters.

\chapter*{Conclusions}
\addcontentsline{toc}{chapter}{Conclusions}

Symbulator is a useful tool for the analysis of linear circuits. It
allows the use of numerical and symbolic values. It can simulate in
direct current, alternating current, the frequency domain and
transient analysis. It gives as answers all the values that interest
the user: voltages at the nodes, and currents, voltage drops and
powers in the elements. It makes it easy to find the Th\'evenin,
Norton and two-port equivalents of a network.

Because of its capabilities, Symbulator can be of great help in the
teaching, learning and practice of Electrical Engineering, in the
area of linear circuit analysis, especially in those circuit
problems that require symbolic answers.

It is to be hoped that the technology developed in Symbulator, and
matured during these two years of polishing, will serve as the basis
for developing more and better products in the future, including
programs for new platforms such as personal computers and portable
computers.

\chapter*{Recommendations}
\addcontentsline{toc}{chapter}{Recommendations}

I do not write these recommendations as the author, but based on my
own experience as a user, and on the testimonials I have received in
writing from the hundreds of students and professionals who have
used Symbulator as a learning and verification tool, during its two
years of development.

I recommend that professors use Symbulator as a teaching instrument
in their circuit theory courses. I believe this simulator has a
place in the classroom, due to its characteristics and portability.
Professors can use it to verify the answer keys of their
examinations before grading them, and as part of their classes.

I consider that the student should learn, in the first place, the
manual techniques of circuit solving, and that in evaluations the
professors should require them to present all calculations and steps
manually. At the same time, I recommend that the student be allowed
to use instruments like Symbulator, as a verification tool in
classes, homework and examinations. I think that having a way to
verify their answers will favor the learning and the self-confidence
of the student. Especially because circuit theory textbooks do not
include answers to all the problems they present, and because some
of the answers they do present are wrong.

\chapter*{References}
\addcontentsline{toc}{chapter}{References}

\begin{itemize}
\item \emph{Engineering Circuit Analysis} (An\'alisis de circuitos
  en ingenier\'ia), Hayt \& Kemmerly, fifth edition (third edition
  in Spanish), McGraw-Hill, 1993.
\item \emph{User's Manual --- TI-89 Calculator}, Texas Instruments,
  Banta Book Group, 1998.
\item \emph{User's Manual --- TI-92 Plus Calculator}, Texas
  Instruments, Banta Book Group, 1999.
\end{itemize}

\appendix

\chapter{A brief history of Symbulator}

The original thesis carried, as its first appendix, a brief history
of Symbulator and a note on the language of the program and its
documentation. When the translators began their work, they moved
that material forward into the body of the book, where you have
already read it: it is section 1.4, \emph{Brief history of
Symbulator}, and its subsection on language. It is not repeated
here; this appendix remains in the plan of the book as a signpost,
so that the numbering of the remaining appendices matches the
original.

\chapter{Symbulator and symbolic simulation}

The knowledge that man has of electricity has been developing for
centuries, from Volta's frogs to quantum transistors. The tools of
electrical engineering analysis have developed as well. Symbolic
simulation, being a very recent invention, is among the newest of
these tools.

In 1998 I came upon the first symbolic circuit simulator I had seen
in my life (and in fact the only one I know of apart from
Symbulator). Its name is Syrup, and it is designed to run in the
environment of the mathematical program Maple~V for personal
computers. Although I never came to master this simulator, thanks to
it I understood that symbolic simulation is a very powerful tool for
learning and design. At that time I was a student in the Circuit
Theory~II course at the Azuero Regional Center of the Technological
University, and I felt it would be very useful to have a symbolic
simulator available in the classroom. I struck up a good friendship
with Joe Riel, the author of Syrup, and he explained to me the
underlying concept of this program, an extremely simple concept in
truth. Unlike numerical simulators, including some I myself had
programmed up to that date, Syrup did not use matrices and modified
nodal analysis. On the contrary, it generated nodal equations,
solved them and presented them symbolically. With this idea in mind,
on March 30, 1999, I undertook the programming of a similar
simulator on my TI-89 calculator, which I initially called SCS (for
Symbolic Circuit Simulator) and later Symbulator. Syrup was an
initial inspiration, but not a model to follow, because my idea of
the perfect simulator for the classroom differed in some respects
from Syrup. Since, during the two years it took me to perfect it, I
had no contact with any other symbolic simulator, I was forced to
invent and develop by my own means the techniques of symbolic
simulation that I present today, and which have rendered excellent
results to the thousands of users that Symbulator has in the whole
world. I have made this clarification because I consider that the
work presented in this book should be taken for what it truly is: a
pioneering contribution to a virgin field of electrical engineering.
Perhaps that is why it earned first place in the IEEE Student Paper
Contest 2000 for Latin America (Region~9).

\chapter{When is it preferable not to use Symbulator?}

An important skill that a Symbulator user must develop is that of
recognizing when the simulation represents a saving of time and when
it does not. This skill manifests itself as a kind of scent, of
intuition, that tells you whether, due to the nature of a problem,
it will be more complicated to solve it with a symbolic simulation
than with manual techniques. It is acquired after much practice.
Only your own experience will provide you with this special
intuition.

Below, I present an example that illustrates a case in which it is
advisable to decline the use of the simulator. It is a circuit that,
at first sight, seems to be a perfect candidate for symbolic
simulation in Symbulator. However, it turns out to be much simpler
to use a single line with the calculator's \code{solve} command to
solve it.

\problem{089}

\emph{Problem: In the circuit shown, it is known that
$\mathrm{abs}(I_1) = 7$~A and $\mathrm{abs}(I_2) = 5$~A. Find $I_1$
and $I_2$.}

Solution: We will simply use the \code{cSolve} command:

\begin{entry}
cSolve(abs(i1)=7 and abs(i2)=5 and i1+i2=10,
       {i1=1+1i,i2=1-1i})
\end{entry}

Note two aspects of this line:

We have written three equations to find two unknowns. This might
seem to be an error, since one should have as many equations as
unknowns. However, when we remember that each complex unknown is
really two unknowns, we understand that we wish to solve for four
unknowns and not two. The third equation can be considered, from
this perspective, as two equations instead of one, since in it the
following two equations are implicitly expressed:
\code{real(i1)+real(i2)=10 and imag(i1)+imag(i2)=0}. So, in the
command line shown above we have virtually four equations with four
unknowns.

We have added an equal sign and a value following the name of each
variable, when listing the unknowns. Instead of saying
\code{\{i1,i2\}}, we have said \code{\{i1=1+1i,i2=1-1i\}}. These
values are known as initial guesses, and their function is to give
the calculator a starting point for the numerical resolution of the
equations. The closer the starting point is to the real value, the
faster the answer will be found. Observe that we have given a
positive imaginary value to $I_1$, and a negative imaginary value to
$I_2$. The reason is simple: we know beforehand that, since the
branch of $I_2$ has a capacitor, its imaginary part will be
negative, while the imaginary part of $I_1$ will be positive, since
there is an inductance in that branch.

We press \key{ENTER}. The answers appear. Depending on the modes of
the calculator, the answers will be shown in rectangular or polar
form. Since we wish to see them in polar form, regardless of the
modes, we use the \code{absang} tool, like this:

\begin{entry}
sq\absang(i1)|ans(1)
\end{entry}

\code{ans(1)} refers to the answers obtained by the previous
command. The number in parentheses indicates which previous answer
we are referring to. We press \key{ENTER}. We obtain 7.∠27.66°.

\begin{entry}
sq\absang(i2)|ans(2)
\end{entry}

We press \key{ENTER}. We obtain 5.∠$-$40.54°.

Note how we have obtained these answers easily, using \code{cSolve}.
Using Symbulator, the effort would have been much greater. That is
why, in this case and in many other similar cases, it is preferable
to decline the use of the simulator, and take another, simpler road.

\chapter{Origin of the problems used}

The problems in this book were, in their majority, taken from the
following book: \emph{An\'alisis de circuitos en ingenier\'ia}, Hayt
\& Kemmerly, fifth edition (third edition in Spanish), McGraw-Hill,
1993. The exceptions have been noted. The original thesis closes
with a table giving, for each of the 89 problems, the page of that
book from which it was taken; the table is preserved in the archived
Spanish original, and its exceptions are worth recording here:
Problem~057 was contributed by Nevin McChesney (USA); Problem~066 by
Jake Adams (USA); Problem~087 was set by a professor of the
Technological University of Panama, as was Problem~088 (see
Chapter~12).
